% Generated draft: requires mathematical and source review.
\documentclass[11pt]{article}
\usepackage[margin=25mm]{geometry}
\usepackage{fontspec}
\setmainfont{Latin Modern Roman}
\usepackage{amsmath,amssymb,mathtools,unicode-math}
\setmathfont{Latin Modern Math}
\usepackage{enumitem,needspace,xurl,hyperref}
\setlength{\parindent}{0pt}
\setlength{\parskip}{5pt}
\setlength{\emergencystretch}{3em}
\newcommand{\N}{\mathbb{N}}
\newcommand{\Z}{\mathbb{Z}}
\newcommand{\Q}{\mathbb{Q}}
\newcommand{\R}{\mathbb{R}}
\newcommand{\C}{\mathbb{C}}
\newcommand{\F}{\mathbb{F}}
\newcommand{\A}{\mathbb{A}}
\newcommand{\PP}{\mathbb{P}}
\newcommand{\E}{\mathbb{E}}
\newcommand{\eps}{\varepsilon}
\newcommand{\dd}{\,\mathrm{d}}
\begin{document}
\section*{UnsolvedMath ID: 11100039}
\textbf{Source problem code:} catalogue entry
\subsection{Definitions and mathematical statement}
Let \(\mathcal C_U^G\) be the equivariant stable homotopy category indexed on the source's \(G\)-universe \(U\). In the usual compact-Lie setting a complete universe contains every irreducible real \(G\)-representation with infinite multiplicity; an incomplete universe need not. With sphere object \(S_U\) and internal function spectrum \(F_U\), define
\[
 DX=F_U(X,S_U),\qquad \eta_X:X\longrightarrow D^2X
\]
using the natural biduality map. Dualizability requires evaluation/coevaluation satisfying the duality identities, equivalently the appropriate canonical maps \(DX\wedge Y\to F_U(X,Y)\) to be equivalences. The source's weaker proposed condition is \(\eta_X\) an equivalence, not merely an unspecified abstract isomorphism with a double dual.

The problem is to understand the stable-homotopy axioms and consequences when the chosen generators are not dualizable, particularly for incomplete-universe equivariant spectra, and find a valid replacement condition. Specifically, are those generators weakly dualizable in this biduality sense? Under the usual orbit-generator interpretation, this asks about \(X=\Sigma^\infty_U(G/H)_+\) and the prescribed suspensions, for the allowed subgroups \(H\). The source does not specify the group class, universe, model, exact generating family or which memoir results the replacement should preserve, so no universal replacement theorem is asserted.

Hovey says the memoir's dualizable-generator assumption poorly handles this setting and credits Peter May with having argued against it. The summary describes later equivariant/global theory as progress with appropriate finiteness notions, but gives no exact theorem establishing biduality for these generators.
\subsection{Short English statement}
Find axioms replacing dualizable generators for incomplete-universe equivariant stable homotopy, and test whether the intended generators satisfy natural double-dual equivalence.
\subsection{Sources}
\begin{itemize}
\item Mark Hovey, \emph{Algebraic Topology Problem List}, ``Axiomatic stable homotopy,'' item 7, surviving mirror \url{https://www-users.cse.umn.edu/~tlawson/hovey/}. Recorded access 2026-07-29; extraction html-ordered-list.
\item Peter May, source attribution for the objection to the memoir axiom. The memoir and later equivariant/global-theory progress have no complete references in this record.
\item Local UnsolvedMath registry, \texttt{data/unsolvedmath/problems.json}, permanent ID \(11100039\), exact source code \texttt{catalogue entry}; supplied partially\_solved/PARTIAL-PROGRESS summary describes conceptual resolution without a single precise theorem.
\end{itemize}
\end{document}
